\section{Five use cases}
\label{sec:usecases}

Each use case is a realistic question, with its application and
configuration files in \cmd{examples/}. Together they use every key of the
two files. Run them from the root of the repository; the outputs below
come from these exact commands.

\begin{center}\small\row
\begin{tabularx}{\textwidth}{@{}>{\raggedright\arraybackslash}p{45mm}X@{}}
\toprule\rowcolor{white}
use case & what it shows \\
\midrule
1.\ sizing the central computer & a zonal platform, 24 vehicle functions; read every output; compare the rules; size for a margin \\
2.\ a fail-operational chassis & redundant task-chains, lock-step cores, a ring with frame replication (FRER), CADU \\
3.\ camera streams and control loops & a tree backbone, three traffic classes, TAS, ATS, frame preemption, link rates \\
4.\ software updates & sizing at the SOP, then updates on frozen hardware; why spreading the load helps then \\
5.\ a hybrid platform & big.LITTLE central computer, a domain controller, shared-memory contention, a custom backbone \\
\bottomrule
\end{tabularx}
\end{center}

% ---------------------------------------------------------------------------
\subsection{Use case 1: how many cores does the central computer need?}
\label{sec:uc1}

\paragraph{The question.} A carmaker moves 24 functions of chassis,
powertrain, ADAS, body and comfort onto a zonal platform: a
high-performance computer (HPC) and five zone controllers. The zone
controllers are chosen (two microcontroller cores at half the speed of the
reference core, 4\,MB for the tasks); the HPC variant is open. How many
cores must it have, and is every function on time?

\begin{figure}[ht]
\centering
\begin{tikzpicture}
\draw[cgrid,line width=1.2pt,rounded corners=18pt,fill=cpaper] (0,0) rectangle (13,5.2);
\node[font=\sffamily\scriptsize,text=cmuted,rotate=90] at (12.6,2.6) {front};
\node[font=\sffamily\scriptsize,text=cmuted,rotate=90] at (0.4,2.6) {rear};
\node[ecuzone] (fl) at (10.8,4.15) {ZC\_FL\cores{2 MCU 50\,\%}};
\node[ecuzone] (fr) at (10.8,1.05) {ZC\_FR\cores{2 MCU 50\,\%}};
\node[ecuzone] (rl) at (2.2,4.15) {ZC\_RL\cores{2 MCU 50\,\%}};
\node[ecuzone] (rr) at (2.2,1.05) {ZC\_RR\cores{2 MCU 50\,\%}};
\node[ecuzone] (c) at (4.9,2.6) {ZC\_C\cores{2 MCU 50\,\%}};
\node[ecucentral,minimum width=24mm] (hpc) at (8.5,2.6) {HPC\cores{2 big, up to 16}};
\node[switch] (sw) at (6.45,2.6) {SW};
\draw[link] (fl) -| (sw); \draw[link] (fr) -| (sw);
\draw[link] (rl) -| (sw); \draw[link] (rr) -| (sw);
\draw[link] (c) -- (sw); \draw[link] (hpc) -- (sw);
\node[font=\sffamily\tiny,text=cmuted,anchor=north] at (6.45,2.25) {1 Gbit/s star};
\end{tikzpicture}
\caption{Use case 1: the platform, seen from above (front on the right).
Each zone controller has 4\,MB and at most 4 cores; the HPC 64\,MB and at
most 16 cores.}
\label{fig:uc1}
\end{figure}

\paragraph{The files.} \cmd{examples/uc1\_central\_sizing/}: the
configuration lists the six ECUs; the HPC starts with only 2 big cores,
so that the tool adds the cores it needs (\key{max\_cores} 16). The
application has 24 task-chains and 73 tasks: brake-by-wire and
steer-by-wire at 1\,ms with 5\,ms deadlines, stability and propulsion
control at 5 and 10\,ms, ADAS perception at 40 to 100\,ms (big cores only,
up to 16\,MB), body and comfort functions at 20\,ms to 1\,s, and three
local chains (two window lifters, the seat heating) bound to one zone
controller each.

\paragraph{Run.}
\begin{shell}
mapping examples/uc1\_central\_sizing/application.json \textbackslash\\
\quad examples/uc1\_central\_sizing/configuration.json -{}-cores
\end{shell}
\Cref{sec:output} reads this output line by line. The answer: \textbf{two
HPC cores suffice} (7 cores in all, the lower bound being 5), every core
is schedulable with a margin of 30\,\% on every WCET, and all 24 functions
meet their deadlines.

\paragraph{Compare the rules.} \cmd{-{}-algorithm all} maps with the four
rules:
\listing{listings/uc1_rules.txt}
CADU saves a core and a third of the network messages here, but one chain
misses its deadline; WFDU spreads the load over 12 cores and also leaves
one chain late. The largest maximum reaction time is about the same for
all rules: it is the chain of the tire pressure monitor, whose tasks run
once a second and whose deadline is 3\,s.

\paragraph{Size for a margin.} A margin of 30\,\% on every WCET is
little at the start of a project. \opt{target-scaling} lowers the
utilization cap until the critical scaling factor reaches the target:
\listing{listings/uc1_sizing.txt}
The cores are the same, but the load is spread so that every WCET may
grow by 68\,\% (target 1.5) or 100\,\% (target 2). The price is paid in
deadlines: with a lower cap, the control tasks of the 1\,ms by-wire chains
no longer fit the zone controller of their sensor and move to the HPC,
which adds two network hops. With the target 1.5 the steer-by-wire chain
is late, with the target 2 the two brake-by-wire chains too; the
\opt{cores} output names them.

% ---------------------------------------------------------------------------
\subsection{Use case 2: a fail-operational chassis}
\label{sec:uc2}

\paragraph{The question.} Steer-by-wire and brake-by-wire must survive
the failure of one ECU or one network link. Each is built as two
redundant task-chains (two \emph{lanes}) with their own sensors and
actuators; their safety monitors must run on lock-step cores; the
messages of the redundant chains are sent twice, over both directions of
an Ethernet ring (IEEE 802.1CB FRER). Is the platform of \cref{fig:uc2}
enough, and what does the redundancy cost on the network?

\begin{figure}[ht]
\centering
\begin{tikzpicture}
\draw[cgrid,line width=1.2pt,rounded corners=18pt,fill=cpaper] (0,0) rectangle (13,5.6);
\node[font=\sffamily\scriptsize,text=cmuted,rotate=90] at (12.6,2.8) {front};
\node[ecuzone] (fl) at (10.6,4.55) {ZC\_FL\cores{2 MCU + 1 lock-step}};
\node[ecuzone] (fr) at (10.6,1.05) {ZC\_FR\cores{2 MCU + 1 lock-step}};
\node[ecuzone] (rl) at (2.4,4.55) {ZC\_RL\cores{2 MCU + 1 lock-step}};
\node[ecuzone] (rr) at (2.4,1.05) {ZC\_RR\cores{2 MCU + 1 lock-step}};
\node[ecucentral] (a) at (5.2,2.8) {HPC\_A\cores{4 big}};
\node[ecucentral] (b) at (7.8,2.8) {HPC\_B\cores{4 big}};
\node[switch] (sa) at (5.2,4.0) {}; \node[switch] (sb) at (7.8,4.0) {};
\node[switch] (sfl) at (8.9,4.55) {}; \node[switch] (sfr) at (8.9,1.05) {};
\node[switch] (srl) at (4.1,4.55) {}; \node[switch] (srr) at (4.1,1.05) {};
\draw[link] (a) -- (sa); \draw[link] (b) -- (sb);
\draw[link] (fl) -- (sfl); \draw[link] (fr) -- (sfr);
\draw[link] (rl) -- (srl); \draw[link] (rr) -- (srr);
% ring in ECU order: A, B, FL, FR, RR, RL, A
\draw[fast,accent] (sa) -- (sb) -- (sfl);
\draw[fast,accent] (sfl) -- (sfr);
\draw[fast,accent] (sfr) -- (srr);
\draw[fast,accent] (srr) -- (srl);
\draw[fast,accent] (srl) -- (sa);
\node[font=\sffamily\tiny,text=accent!80!black] at (6.5,0.55) {ring of six switches, 1 Gbit/s, FRER};
\end{tikzpicture}
\caption{Use case 2: two central computers and four zone controllers with
lock-step cores, one switch per ECU, the switches in a ring in the order
of the configuration file.}
\label{fig:uc2}
\end{figure}

\paragraph{The files.} \cmd{examples/uc2\_redundant\_chassis/}. The
configuration sets a cap of 0.7 and a WCET margin of 10\,\% (safety
functions), a ring with FRER, two CBS classes (messages of tasks up to
5\,ms in the first one, with 50\,\% of the port, the others with 25\,\%),
and the CADU rule with replicas on different ECUs. In the application,
\key{replica\_group} 1 joins the two steering lanes and group 2 the two
braking lanes; the monitors have \key{"core\_types": ["lockstep"]} and the
ECUs they may use (\key{ecus}); the airbag is a local chain of the front
right zone on lock-step cores; the state estimator and the torque
vectoring of the vehicle dynamics share a core (\key{colocate\_with\_next});
the diagnostics may only run on the two HPCs (a chain-level \key{ecus}).

\paragraph{Run.}
\begin{shell}
mapping examples/uc2\_redundant\_chassis/application.json \textbackslash\\
\quad examples/uc2\_redundant\_chassis/configuration.json -{}-cores
\end{shell}
\listing[\fontsize{4.8pt}{5.8pt}\selectfont]{listings/uc2_chains.txt}
CADU keeps each lane in its own zones: the steering lanes never leave
\cmd{ZC\_FL} and \cmd{ZC\_FR} (one hop between two cores of the zone
controller, through its shared memory), so the replicas are on different
ECUs by construction, and only the braking lanes cross the ring: two
replicated flows. Every chain is on time with 8 cores.

\paragraph{The price of fewer messages.}
\listing{listings/uc2_rules.txt}
BFDU and FFDU need only 6 cores, but they send the control tasks of both
lanes to a central computer: twice as many messages, more than twice as
many replicated flows, and network bounds six times larger. With
\opt{separate-cores} the replicas may share an ECU (not a core): try it
to see which constraint costs what. A tree topology is refused with
\opt{frer}: a tree has no second route.

% ---------------------------------------------------------------------------
\subsection{Use case 3: camera streams and control loops on one backbone}
\label{sec:uc3}

\paragraph{The question.} Cameras (up to 120\,KB per frame), radars and a
lidar send their data from the zone controllers to the central computer,
while 1\,ms control loops cross the same switches between zones. Which
traffic shaping keeps the control loops on time, and what happens with a
slower link?

\begin{figure}[ht]
\centering
\begin{tikzpicture}
\node[switch,minimum width=10mm] (root) at (6.5,2.4) {SW};
\node[ecucentral,minimum width=26mm] (hpc) at (6.5,3.75) {HPC\cores{8 big 100\,\%, 256\,MB}};
\draw[fast,brand] (hpc) -- node[right,font=\sffamily\tiny,text=brand]{10 Gbit/s} (root);
\foreach \n/\x/\s [count=\k] in {ZC\_F/1.5/front camera\, radar\, lidar,ZC\_L/4.8/left camera\, corner radar,ZC\_R/8.2/right camera,ZC\_B/11.5/rear camera} {
	\node[switch] (s\k) at (\x,1.0) {};
	\draw[fast] (root) -- node[pos=0.55,font=\sffamily\tiny,text=cmuted,fill=white,inner sep=1pt]{2.5 Gbit/s} (s\k);
	\node[ecuzone,minimum width=24mm] (z\k) at (\x,-0.3) {\n\cores{3 MCU 50\,\%}};
	\draw[link] (s\k) -- (z\k);
	\node[font=\sffamily\tiny,text=accent!80!black,anchor=north,text width=26mm,align=center] at (\x,-1.0) {\s};
}
\end{tikzpicture}
\caption{Use case 3: a tree rooted at the switch of the HPC; the zone
links run at 1\,Gbit/s, the switches are linked at 2.5\,Gbit/s and the
HPC at 10\,Gbit/s.}
\label{fig:uc3}
\end{figure}

\paragraph{The files.} \cmd{examples/uc3\_tsn\_streams/}. The network
model is a \key{tree} with \key{tree\_root} \cmd{HPC},
\key{backbone\_mbps} 2500 and \key{ecu\_link\_mbps} 10000 for the HPC;
three traffic classes (producers up to 5\,ms, up to 40\,ms, the rest), TAS
with a cycle of 500\us{} and windows of 20, 40 and 20\,\%, frame
preemption, and a design limit of 80\,\% per port. The application has 12
chains: perception streams (40 to 100\,ms, messages of 3 to 120\,KB) and
control loops (1 to 20\,ms, messages of 8 to 32 bytes).

\paragraph{Run, then change the network.}
\begin{shell}
mapping examples/uc3\_tsn\_streams/application.json \textbackslash\\
\quad examples/uc3\_tsn\_streams/configuration.json [option]
\end{shell}
\listing{listings/uc3_network.txt}
With the file as it is, every chain is on time; the largest delay bound
(14\,ms) is the one of a camera frame of 120\,KB, sent through the TAS
window of its class. ATS (\opt{ats}) removes the waiting for the gates:
the largest bound falls to 2.9\,ms. With one class only
(\opt{classes 1}) the messages of the 1\,ms loops wait behind the camera
frames: the steering and braking loops miss their deadline, and so does
the front camera chain. With 100\,Mbit/s links
(\opt{link 100}) the ports carry 21\,\% of LET traffic and five chains are
late; ATS saves four of them.

% ---------------------------------------------------------------------------
\subsection{Use case 4: software updates on frozen hardware}
\label{sec:uc4}

\paragraph{The question.} The hardware of a car is fixed at the start of
production (SOP); its software keeps growing with updates for years. The
platform of use case~1 is sized at the SOP with a margin; then the HPC
variant is frozen. Do two later updates still fit, and how should they be
mapped?

\begin{figure}[ht]
\centering
\begin{tikzpicture}[node distance=6mm]
\node[box,minimum width=36mm,text width=34mm] (sop) {\textbf{SOP}\\24 functions\\\scriptsize sized with target 1.5};
\node[box,right=14mm of sop,minimum width=36mm,text width=34mm,fill=white] (u1) {\textbf{update 1}\\+4 functions\\\scriptsize energy, lane centering\dots};
\node[box,right=14mm of u1,minimum width=36mm,text width=34mm,fill=white] (u2) {\textbf{update 2}\\+4 functions\\\scriptsize highway pilot};
\draw[arrow] (sop) -- (u1); \draw[arrow] (u1) -- (u2);
\node[draw=accent,fill=accentlight,rounded corners=2pt,below=5mm of u1,
	minimum width=104mm,font=\sffamily\scriptsize,align=center] (hw)
	{frozen hardware (\key{fixed\_hardware}): HPC with 3 big cores, each zone controller with its 2 cores};
\end{tikzpicture}
\caption{Use case 4: one platform, three versions of the software.}
\label{fig:uc4}
\end{figure}

\paragraph{The files.} \cmd{examples/uc4\_software\_updates/}:
\cmd{application\_sop.json} (the functions of use case~1),
\cmd{application\_update1.json} and \cmd{application\_update2.json} (four
more functions each, among them a highway pilot with 32\,MB of
perception); \cmd{configuration\_sop.json} (use case~1 with
\key{"target\_scaling": 1.5} and 512\,MB on the HPC) and
\cmd{configuration\_frozen.json} (the hardware chosen: 3 big cores on the
HPC, at most the initial cores everywhere, \key{"fixed\_hardware": true}).

\paragraph{Size at the SOP.}
\begin{shell}
mapping examples/uc4\_software\_updates/application\_sop.json \textbackslash\\
\quad examples/uc4\_software\_updates/configuration\_sop.json
\end{shell}
\listing{listings/uc4_sop.txt}
The SOP needs 2 HPC cores for a scaling factor of 1.68; the carmaker buys
a third one for the updates.

\paragraph{Map the updates on the frozen platform.} Each version is
mapped with the four rules on \cmd{configuration\_frozen.json}
(\cmd{-{}-algorithm all}); the table gives the task-chains on time:
\listing{listings/uc4_updates.txt}
Both updates fit the frozen hardware with every rule but CADU, which
cannot place a sensor task bound to \cmd{ZC\_FL} in update~2: no core of
that zone controller has room left, and the hardware is fixed. But the
rules that minimize the cores (BFDU, FFDU) now leave chains late:
packing the new functions moves the steering control off its zone
controller. \textbf{On frozen hardware the cores are already paid for}:
WFDU, which spreads the load over every core, keeps all chains of update~1
on time and 31 of 32 of update~2. The allocation objective changes after
the SOP.

% ---------------------------------------------------------------------------
\subsection{Use case 5: a hybrid platform with a big.LITTLE central computer}
\label{sec:uc5}

\paragraph{The question.} A hybrid platform keeps an ADAS domain
controller next to a central computer whose cores are of two types: 4
big and 4 little (at 40\,\% of the speed). The cores of each ECU share its
memory. Where does each function go, and how much does memory contention
cost?

\begin{figure}[ht]
\centering
\begin{tikzpicture}
\node[switch,minimum width=16mm] (sc) at (6.5,2.6) {SW\_CENTRAL};
\node[switch,minimum width=14mm] (sf) at (10.2,1.0) {SW\_FRONT};
\node[switch,minimum width=14mm] (sr) at (2.8,1.0) {SW\_REAR};
\node[ecucentral,minimum width=30mm] (hpc) at (6.5,3.9) {HPC\cores{4 big 100\,\% + 4 little 40\,\%}};
\draw[fast,brand] (hpc) -- node[right,font=\sffamily\tiny,text=brand]{10G} (sc);
\draw[fast] (sc) -- node[above right,font=\sffamily\tiny,text=cmuted]{10G} (sf);
\draw[fast] (sc) -- node[above left,font=\sffamily\tiny,text=cmuted]{2.5G} (sr);
\draw[link] (sf) -- node[below,font=\sffamily\tiny,text=cmuted]{1G} (sr);
\node[ecudomain,minimum width=24mm] (adas) at (12.0,2.6) {ADAS\_DC\cores{4 big 120\,\%}};
\draw[fast,teal] (adas) -- node[right,font=\sffamily\tiny,text=teal]{10G} (sf);
\node[ecuzone] (fl) at (8.9,-0.4) {ZC\_FL\cores{2 MCU 50\,\%}};
\node[ecuzone] (fr) at (11.6,-0.4) {ZC\_FR\cores{2 MCU 50\,\%}};
\node[ecuzone] (rl) at (1.4,-0.4) {ZC\_RL\cores{2 MCU 50\,\%}};
\node[ecuzone] (rr) at (4.2,-0.4) {ZC\_RR\cores{2 MCU 50\,\%}};
\draw[link] (fl) -- (sf); \draw[link] (fr) -- (sf);
\draw[link] (rl) -- (sr); \draw[link] (rr) -- (sr);
\end{tikzpicture}
\caption{Use case 5: a custom backbone of three switches; the links of
the HPC and of the ADAS controller run at 10\,Gbit/s.}
\label{fig:uc5}
\end{figure}

\paragraph{The files.} \cmd{examples/uc5\_hybrid\_platform/}. The
configuration has every section: the platform (a \key{central}, a
\key{domain} and four \key{zone} ECUs, a 120\,\% ADAS controller), the
interference model (40\,ns per memory request), a \key{custom} topology
(\cref{sec:conf-network}), ATS with two classes, and CADU. In the
application the neural detector gains 30\,\% on big cores
(\key{"wcet\_factor\_percent": \{"big": 70\}}) and may only run on the ADAS
controller; the control laws lose less than the speed ratio on little
cores (\key{\{"little": 80\}}); the battery and comfort strategies may
only use little cores; every task gives its memory requests per job
(\key{memory\_accesses}).

\paragraph{Run.}
\begin{shell}
mapping examples/uc5\_hybrid\_platform/application.json \textbackslash\\
\quad examples/uc5\_hybrid\_platform/configuration.json -{}-cores
\end{shell}
\listing{listings/uc5_cores.txt}
The perception runs on the ADAS controller (core c8), the planning on a
big HPC core (c0), the strategies of battery, climate and body and the
diagnostics on one little HPC core (c4); the control loops stay in the
zone controllers. All chains are on time.

\paragraph{The cost of contention.} \opt{memory-latency} overrides the
interference model:
\listing{listings/uc5_contention.txt}
Without contention 6 cores suffice. At 40\,ns per request, the WCETs grow
with the number of cores competing for the memory of each ECU, and a
seventh core is needed; at 100\,ns an eighth one, and the application
load grows by 10\,\%. The tool finds the number of competing cores of
each ECU as a fixed point (more cores, more contention, longer WCETs).
